\section{Conclusion}
\label{sec:conclusion}

We presented \method{}, a failure-triggered in-loop memory system for ReAct agents operating in recurring-instance settings. Unlike episode-boundary approaches, \method{} monitors execution in real time, detects failures, and retrieves relevant past experiences to provide mid-episode course correction. The formal ScienceWorld Step~4 protocol shows both the promise and limits of in-loop memory: online \method{} reaches 41.44\% success rate and the highest raw score among non-ExpeL memory settings (64.25) on 111 episodes, improving over offline memory (39.64\%, 48.83), the ReAct baseline (37.84\%, 47.21), and corrected Reflexion (36.94\%, 49.39), while ExpeL reaches the highest success rate after a full insight-collection epoch (49.55\%, 57.78). Across the broader benchmark suite, our results suggest that in-loop memory is most useful when online correction matters, while episode-level methods remain competitive or stronger when full-trajectory strategy extraction is available.

Experiments also show the highest success rate on ALFWorld, while ExpeL achieves the best success-rate results on ScienceWorld, WebShop, HotPotQA, and ToolBench. The ALFWorld/WebShop ablations indicate that the timing advantage of in-loop retrieval is conditional rather than universal: in-loop retrieval is directionally stronger on ALFWorld but slightly weaker than episode-level injection on WebShop. The structured failure-recovery format still outperforms abstract reflections by +3.7\% on ALFWorld ($p{=}0.056$, borderline), while the advantage over success trajectories is directional but not significant. A cross-environment memory transfer experiment across three benchmarks reveals that shared memory effectiveness depends on task homogeneity: on WebShop (homogeneous tasks), shared memory achieves Epoch~2 performance in a single epoch; on HotPotQA, moderate transfer occurs (+1.3\% EM); on ToolBench (heterogeneous APIs), shared memory causes negative transfer ($-$19.9\%), highlighting the need for selective sharing based on task similarity.

\paragraph{Limitations.} First, our main evaluation studies online adaptation on recurring task instances (the same environments across epochs). The cross-environment transfer experiment (\S\ref{sec:cross_env}) shows that transfer effectiveness depends on task homogeneity---positive on WebShop, moderate on HotPotQA, and negative on ToolBench---indicating that shared memory requires careful deployment. Second, the implicit failure detector relies on an LLM judge, but our precision audit in the final protocol is centered on WebShop; the formal ScienceWorld Step~4 result instead uses rule-based explicit detection. We also did not measure recall---the fraction of true failures that the detector catches---which would require labeling all action-observation pairs, not just flagged ones. Third, different LLM backbones are used across benchmarks (Qwen3.5-9B for ALFWorld/WebShop, DeepSeek-V3.2 for ScienceWorld), which limits cross-benchmark comparisons to trends rather than direct comparisons. Fourth, the formal ScienceWorld Step~4 comparison is a dedicated protocol rather than a clean timing-only ablation; it should also not be interpreted as a direct leaderboard comparison against ScienceWorld papers that use different task subsets, backbones, multi-trial adaptation, or trained components. Finally, on ScienceWorld, failure memory improves performance but does not solve the persistent hard cases, especially \texttt{mendelian-genetics-unknown-plant}, which remains unsolved in all Step~4 settings.

\paragraph{Future Work.} Promising directions include similarity-aware cross-environment memory sharing (e.g., clustering environments by task similarity and sharing memory only within clusters to avoid negative transfer), memory consolidation to merge redundant entries and distill higher-level strategies, and adaptive retrieval that adjusts the failure detection threshold based on task difficulty.
